\documentclass{article}
\usepackage[margin=1in, a4paper]{geometry}
\usepackage{amsmath, amssymb, amsfonts}
\usepackage{graphicx}
\usepackage{xcolor}
\usepackage{listings}
\usepackage{booktabs}
\usepackage[utf8]{inputenc}
\usepackage[T1]{fontenc}
\usepackage{hyperref}
\hypersetup{colorlinks=true, urlcolor=blue, linkcolor=blue}
\usepackage{enumitem}
\usepackage{float}

\definecolor{codegreen}{rgb}{0,0.6,0}
\definecolor{codegray}{rgb}{0.5,0.5,0.5}
\definecolor{codepurple}{rgb}{0.58,0,0.82}
\definecolor{backcolour}{rgb}{0.99,0.99,0.99}
% Professional color palette for academic publishing
\definecolor{myblue}{cmyk}{0.8,0.4,0,0.1}
\definecolor{mygreen}{cmyk}{0.7,0,0.5,0.2}
\definecolor{myorange}{cmyk}{0,0.5,0.8,0.1}
\definecolor{mygray}{cmyk}{0,0,0,0.4}
\definecolor{arrowcolor}{cmyk}{0.9,0.5,0,0.3}
\definecolor{nodecolor}{cmyk}{0.6,0.2,0,0.05}
\definecolor{framecolor}{cmyk}{0.4,0.1,0,0.1}

\lstdefinestyle{mystyle}{
    backgroundcolor=\color{backcolour},   
    commentstyle=\color{codegreen},
    keywordstyle=\color{magenta},
    numberstyle=\tiny\color{codegray},
    stringstyle=\color{codepurple},
    basicstyle=\ttfamily\footnotesize,
    breakatwhitespace=false,         
    breaklines=true,                 
    captionpos=b,                    
    keepspaces=true,                 
    numbers=left,                    
    numbersep=5pt,                  
    showspaces=false,                
    showstringspaces=false,
    showtabs=false,                  
    tabsize=2
}
\lstset{style=mystyle}

\title{The Container Stacking Rules Problem}
\author{The Logistics \& Supply Chain Problem-Solving Playbook}
\date{}

\begin{document}

\maketitle

\section{The Container Stacking Rules Problem}

Container stacking rules form the backbone of terminal efficiency, directly impacting the number of unproductive container movements known as reshuffles[1][3]. The challenge lies in establishing intelligent stacking policies that minimize these costly repositioning operations while maintaining accessibility to containers based on their planned departure times. This problem becomes increasingly complex as yard utilization increases, with research showing that solutions differ significantly between low-loaded yard-bays (fewer than 15 containers) and high-loaded configurations[2].

\subsection{The Scenario: Optimizing Terminal Yard Operations Through Intelligent Stacking}

Maritime container terminals face a perpetual challenge: how to stack incoming containers in yard bays to minimize the number of reshuffles required when containers need to be retrieved for loading onto outbound vessels. Each reshuffle represents lost productivity, increased fuel consumption, and potential delays in vessel departure schedules.

Consider a typical yard bay with 6 rows and 4-5 tiers of stacking height[2]. Containers arrive with known departure times, destinations, and priorities. The terminal must decide not only where to place each incoming container but also establish rules that can be applied consistently by yard crane operators. The stacking rules must balance multiple objectives: minimizing reshuffles, maintaining efficient space utilization, and ensuring timely access to containers based on their departure schedules.

Two primary stacking strategies emerge from practice: category stacking, where containers of similar characteristics are grouped together, and residence time stacking, where containers are positioned based on their planned departure times[3]. The effectiveness of these strategies depends heavily on the specific operational context and the trade-offs between stacking efficiency and retrieval accessibility.

\subsection{Tier 1: The Pen \& Paper Method (Mathematical Formulation)}

The container stacking rules problem can be formulated as an integer programming model that determines optimal stacking positions while minimizing the expected number of reshuffles.

\textbf{Sets and Parameters:}
\begin{align}
S &= \{1, 2, \ldots, s\} \quad \text{set of stack positions} \\
T &= \{1, 2, \ldots, t\} \quad \text{set of tier levels} \\
C &= \{1, 2, \ldots, n\} \quad \text{set of containers} \\
d_i &= \text{departure time of container } i \\
w_i &= \text{weight of container } i \\
h_{st} &= \text{height capacity at stack } s, \text{ tier } t \\
M &= \text{large constant for logical constraints}
\end{align}

\textbf{Decision Variables:}
\begin{align}
x_{ist} &= \begin{cases} 
1 & \text{if container } i \text{ is placed at stack } s, \text{ tier } t \\
0 & \text{otherwise}
\end{cases} \\
y_{ijst} &= \begin{cases}
1 & \text{if container } i \text{ blocks container } j \text{ at stack } s, \text{ tier } t \\
0 & \text{otherwise}
\end{cases} \\
r_{ij} &= \begin{cases}
1 & \text{if container } j \text{ requires a reshuffle due to container } i \\
0 & \text{otherwise}
\end{cases}
\end{align}

\textbf{Objective Function:}
\[
\min Z = \sum_{i \in C} \sum_{j \in C} r_{ij}
\]

\textbf{Constraints:}
\begin{align}
\sum_{s \in S} \sum_{t \in T} x_{ist} &= 1 \quad \forall i \in C \quad \text{(each container assigned once)} \\
\sum_{i \in C} x_{ist} &\leq 1 \quad \forall s \in S, t \in T \quad \text{(position capacity)} \\
\sum_{i \in C} x_{is,t'} &\geq x_{jst} \quad \forall j \in C, s \in S, t \in T, t' < t \quad \text{(stacking order)} \\
y_{ijst} &\geq x_{ist} + x_{js,t'} - 1 \quad \forall i,j \in C, s \in S, t > t' \quad \text{(blocking detection)} \\
r_{ij} &\geq y_{ijst} \cdot (d_j - d_i)/M \quad \forall i,j \in C, s \in S, t > t' \quad \text{(reshuffle identification)} \\
x_{ist} &\in \{0,1\} \quad \forall i \in C, s \in S, t \in T
\end{align}

\begin{figure}[htbp]
\centering
\includegraphics[width=\textwidth]{../figures-png/line2.fig1.png}
\caption{Container stacking problem visualization showing blocking relationships and reshuffle requirements}
\end{figure}

\textbf{Concrete Example:} Consider 5 containers with departure times: C1(d=3), C2(d=1), C3(d=5), C4(d=2), C5(d=6). With 3 stacks and 2 tiers maximum, the optimal solution places C2 at the bottom of stack 1, C4 at the bottom of stack 2, C1 on top of C2, C3 at the bottom of stack 3, and C5 on top of C3. This configuration requires exactly 1 reshuffle when C2 (earliest departure) needs to be retrieved, as C1 must be moved first.

\subsection{Tier 2: The Classic Heuristic (Python Implementation)}

The constructive heuristic approach builds stacking solutions incrementally by applying priority rules that consider both departure times and stack configurations. This method processes containers in arrival order and assigns each to the best available position based on a multi-criteria evaluation function.

\textbf{Algorithm: Priority-Based Constructive Heuristic}

The algorithm evaluates each potential stacking position using weighted criteria including reshuffle potential, stack height balance, and departure time compatibility.

\lstinputlisting[language=Python]{.././codes-python/line2.py1.py}

\begin{figure}[htbp]
\centering
\includegraphics[width=\textwidth]{../figures-png/line2.fig2.png}
\caption{Enhanced Priority-Based Constructive Heuristic with Node-Based Flow Structure}
\end{figure}

\textbf{Time Complexity Analysis:} The algorithm has a time complexity of O(n × s × h), where n is the number of containers, s is the number of stacks, and h is the average height of stacks. For each container, we evaluate all available stack positions, making this approach efficient for real-time terminal operations.

\textbf{Concrete Example:} Processing containers C1(d=3), C2(d=1), C3(d=5), C4(d=2), C5(d=6) with 3 stacks and 2 tiers maximum. The algorithm assigns: C1 to Stack 1 Tier 0 (score=2.0), C2 to Stack 2 Tier 0 (score=2.0), C3 to Stack 3 Tier 0 (score=2.0), C4 to Stack 2 Tier 1 (score=12.0 due to departure conflict with C2), C5 to Stack 3 Tier 1 (score=2.0). Total expected reshuffles: 1 (C4 blocks C2).

\subsection{Tier 3: The Advanced Algorithm (Metaheuristic Implementation)}

The Genetic Algorithm approach evolves populations of stacking rule configurations to discover near-optimal solutions for complex container stacking scenarios. This metaheuristic excels at exploring the large solution space while avoiding local optima that trap traditional optimization methods[4].

\textbf{Genetic Algorithm for Container Stacking Rules}

The algorithm represents stacking rules as chromosomes encoding position preferences and uses evolutionary operators to improve solution quality over generations.

\lstinputlisting[language=Python]{.././codes-python/line2.py2.py}

\begin{figure}[htbp]
\centering
\includegraphics[width=\textwidth]{../figures-png/line2.fig3.png}
\caption{Enhanced Genetic Algorithm with Node-Based Flow Structure and Professional Chromosome Encoding for Container Stacking Optimization}
\end{figure}

\textbf{Parameter Tuning:} Optimal parameters include population size of 100, mutation rate of 0.1, tournament size of 3, and 500 generations. These parameters balance exploration and exploitation while maintaining computational efficiency.

\textbf{Concrete Example:} For 5 containers with 3 stacks, the genetic algorithm evolves from random initial solutions with 3-4 reshuffles to optimized configurations achieving 0-1 reshuffles. A typical evolved chromosome [2,0,1,2,0] assigns containers to specific stacks, resulting in a stack configuration that minimizes blocking relationships and achieves the theoretical minimum reshuffle count.

\subsection{Tier 4: The AI/ML/RL Augmentation Method}

Reinforcement Learning transforms container stacking from static rule application to dynamic policy learning, where an RL agent discovers optimal stacking strategies through interaction with the terminal environment. The agent learns to balance immediate stacking decisions with long-term reshuffle minimization objectives.

\textbf{Reinforcement Learning Formulation:}

\textbf{State Space:} \(S = (H, D, W, T)\) where:
\begin{itemize}
\item \(H\): Current height configuration of all stacks
\item \(D\): Departure time vector of containers in current stacks
\item \(W\): Weight distribution across stacks
\item \(T\): Time until next departure
\end{itemize}

\textbf{Action Space:} \(A = \{a_{st} : s \in \{1,\ldots,S\}, t \in \{1,\ldots,T\}\}\) representing placement of arriving container at stack \(s\), tier \(t\).

\textbf{Reward Function:}
\[
R(s,a,s') = -\alpha \cdot \text{reshuffles}(s,a,s') - \beta \cdot \text{height\_penalty}(s') + \gamma \cdot \text{efficiency\_bonus}(s')
\]

\lstinputlisting[language=Python]{.././codes-python/line2.py3.py}

\begin{figure}[htbp]
\centering
\includegraphics[width=\textwidth]{../figures-png/line2.fig4.png}
\caption{Enhanced Deep Q-Network Learning Loop with Node-Based Structure and Professional State Representation for Container Stacking}
\end{figure}

\textbf{Training Protocol:} The agent trains over 10,000 episodes, gradually reducing exploration from \(\epsilon = 1.0\) to \(\epsilon = 0.01\). Target network updates occur every 100 episodes to stabilize learning. The replay buffer maintains 10,000 experiences to break temporal correlations in training data.

\textbf{Concrete Example:} After training on historical terminal data, the DQN agent achieves 15\% fewer reshuffles compared to rule-based heuristics. For a test scenario with 20 containers and 4 stacks, the trained agent produces stacking decisions that result in only 2 reshuffles versus 4 reshuffles from traditional first-available placement rules. The learned policy demonstrates sophisticated behaviors like deliberately leaving gaps in stacks to accommodate expected short-residence containers.

\subsection{Tier 8: The Value-Aligned \& Ethical Framework}

The ethical container stacking framework addresses fairness, transparency, and accountability in automated terminal operations. As stacking decisions directly impact vessel schedules, labor allocation, and customer service levels, the system must incorporate value alignment to ensure equitable treatment of all stakeholders while maintaining operational efficiency.

\textbf{Ethical Constraint Architecture}

The value-aligned system implements multi-stakeholder fairness constraints alongside traditional optimization objectives, ensuring that stacking decisions consider broader ethical implications beyond pure efficiency metrics.

\textbf{Stakeholder Value Identification:}
\begin{itemize}
\item \textbf{Shipping Lines:} Fair allocation of premium stacking positions, transparent priority handling
\item \textbf{Terminal Workers:} Ergonomic crane operation patterns, equitable workload distribution
\item \textbf{Environmental:} Minimized fuel consumption, reduced carbon footprint per container move
\item \textbf{Port Community:} Noise reduction during night operations, air quality considerations
\item \textbf{Customers:} Predictable delivery schedules, non-discriminatory service levels
\end{itemize}

\lstinputlisting[language=Python]{.././codes-python/line2.py4.py}

\begin{figure}[htbp]
\centering
\includegraphics[width=\textwidth]{../figures-png/line2.fig5.png}
\caption{Ethical container stacking framework with multi-stakeholder value evaluation}
\end{figure}

\textbf{Algorithmic Accountability Mechanisms:}

The framework implements explainable AI principles through decision logging, stakeholder impact assessment, and regular bias auditing. Each stacking decision generates a transparent explanation that stakeholders can understand and challenge if necessary.

\textbf{Constitutional AI Integration:} The system incorporates constitutional AI principles by defining a ``constitution'' of terminal values that cannot be violated regardless of efficiency gains. These include non-discrimination clauses, environmental protection minimums, and worker safety requirements.

\textbf{Concrete Example:} When processing a batch of 50 containers from different shipping lines, the ethical framework detects that purely efficiency-optimized stacking would systematically disadvantage smaller customers by placing their containers in less accessible positions. The system applies fairness constraints to ensure service level equity, resulting in a 3\% efficiency decrease but eliminating discriminatory treatment patterns. The decision audit trail shows transparent reasoning: ``Efficiency: 87.2, Fairness: 95.1, Environmental: 78.9 - Solution selected to maintain fairness standards while preserving 95\% of optimal efficiency.''

\section{Practice Questions}

\textbf{Question 1 (Tier 1):} Formulate an integer programming model for a container stacking problem with 4 containers having departure times [2, 1, 4, 3], 2 stacks, and maximum 2 tiers per stack. Define all sets, parameters, decision variables, objective function, and constraints. Calculate the optimal solution and total number of reshuffles required.

\textbf{Question 2 (Tier 2):} Implement a priority-based constructive heuristic for the following scenario: 6 containers with departure times [1, 3, 2, 5, 4, 6] and weights [20, 15, 25, 10, 30, 18] tons, 3 stacks, maximum 2 tiers. Use scoring criteria: departure time violations (+10), stack height imbalance (+2), weight instability (+5). Show the step-by-step assignment process and final reshuffle count.

\textbf{Question 3 (Tier 3):} Design a genetic algorithm for container stacking with population size 20, 50 generations, crossover rate 0.8, mutation rate 0.1. Given 5 containers with departure times [2, 1, 4, 3, 5], 2 stacks, 3 tiers maximum, show the chromosome encoding, fitness function calculation, and evolution of the best solution over 5 generations.

\textbf{Question 4 (Tier 4):} Develop a reinforcement learning formulation for dynamic container stacking. Define the state space (include stack heights, departure times, container weights), action space (stack selection), and reward function (reshuffle penalties, efficiency bonuses). For a simple scenario with 3 containers arriving sequentially with departure times [3, 1, 2], show the Q-value updates for the first 3 learning steps assuming initial Q-values are zero and learning rate $\alpha = 0.1$.

\textbf{Question 8 (Tier 8):} Design an ethical evaluation framework for container stacking decisions involving 3 shipping companies: Premium Corp (20 containers), Standard Line (15 containers), and Budget Shipping (10 containers). Define fairness constraints to prevent discrimination, environmental scoring to minimize fuel consumption, and worker welfare considerations for crane operator safety. If a purely efficient solution gives Premium Corp 90\% optimal positions, Standard Line 60\%, and Budget Shipping 40\%, calculate the fairness violation score and propose an alternative solution that maintains 85\% overall efficiency while ensuring equitable treatment (maximum 20\% service level variance between companies).

\end{document}
